\section{Related Work}\label{sec:relatedwork}

\textbf{Formulaic alphas.} The search space of formulaic alphas is enormous, due to the large amount of possible operators and features to choose from. To our best knowledge, all notable former work uses genetic programming to explore this huge search space. \cite{huataigp} augmented the gplearn library with formulaic-alpha-specific time-series operators, upon which an alpha-mining framework is built. \cite{huataigp2} further improved the framework to also mine alphas with non-linear relations with the returns by using mutual information as the fitness measure. \cite{autoalpha} used mutual IC to filter out alphas that are too similar to existing ones, improving the diversity of resulting alpha sets. PCA is carried out on the alpha values for reducing the algorithmic complexity of computing the mutual ICs, and various other tricks are also applied to aid the evolution process. AlphaEvolve \cite{alphaevolve} evolves new alphas upon existing ones. It allows combinations of much more complex operations (for example matrix-wise computations), and uses computation graphs instead of trees to represent the alphas. This leads to more sophisticated alphas and better prediction accuracy, although at the risk of lowering the alphas' interpretability. Mutual IC is also used as a measure of alpha synergy in this work.

\textbf{Machine learning-based alphas.} The development of deep learning in recent years has brought about various new ideas on how to accurately model stock trends. Early work on stock trend forecasting treats the movement of each stock as a separate time series, and applies time series models like LSTM \cite{lstm} or Transformer \cite{transformer} to the data. Specific network structures catered to stock forecasting like the SFM \cite{sfm} which uses a DFT-like mechanism have also been developed. Recently, research has also been conducted on methods to integrate non-standard data with the time series. REST \cite{rest} fuses multi-granular time series data together with historical event data to model the market as a whole. HIST \cite{hist} utilizes concept graphs on top of the regular time series data to model shared commonness between future trends of various stock groups. One specific type of machine learning-based model is also worth mentioning. Decision tree models, notably XGBoost \cite{xgb}, LightGBM \cite{lgbm}, etc., are often considered interpretable, and they could also achieve relatively good performance on stock trend forecasting tasks. However, whether a decision tree with extremely complex structure is considered ``interpretable'' is at least questionable. When these tree models are applied to raw stock data, the high dimensionality of input only exacerbates the aforementioned problem. Our formulaic alphas use operators that apply to the input data in a more structured manner, making them more easily interpretable by curious investors.

\textbf{Symbolic regression.} Symbolic regression (SR) concerns the problem of discovering relations between variables represented in closed-form mathematical formulas. SR problems are different from our problem settings that there always exists a ``groundtruth`` formula that precisely describes the data points in an SR problem, while stock market trends are far too complex to be expressed in the space of formulaic alphas. Nevertheless, there remain similarities between the two fields since similar techniques can be used for the expression generator and the optimization procedure. \cite{srnn} suggested using a custom neural network whose activation functions are symbolic operators to solve the SR problem. \cite{dsr} proposed a novel symbolic regression framework based on an autoregressive expression generator. The generator is optimized using an augmented version of the policy gradient algorithm that values the top performance of the agent more than the average. \cite{dsrgp} developed a method similar to \cite{dsr}, but also introduced GP into the optimization loop, seeding the GP population with RL outputs. \cite{symbolicgpt} applied the language model pretraining scheme to symbolic regression, training a generative autoregressive ``language model'' of expressions on a large dataset of synthetic expressions.

\textbf{Discussions.} Although the term ``formulaic alpha'' is often tied down to investing, the concept of simple and interpretable formulaic predictors that could be combined into more expressive models is not limited to quantitative trading scenarios. Our framework can be adapted to solve other time-series forecasting problems, for example, energy consumption prediction \cite{disc1}, anomaly detection \cite{disc2}, biomedical settings \cite{disc3}, etc. In addition, we chose the linear combination model in this paper for its simplicity. Meanwhile, in theory, other types of interpretable combination models, for example, decision trees can also be integrated into our framework. In that sense, providing these combination models with these features expressed in relatively straightforward formulas might help provide investigators with more insights into how the models come to the final results.
